\documentclass[11pt]{amsart}

\usepackage[margin=1in]{geometry}
\usepackage{amsmath,amssymb,amsthm,mathtools}
\usepackage{enumitem}

\newtheorem{theorem}{Theorem}
\newtheorem{proposition}[theorem]{Proposition}
\newtheorem{lemma}[theorem]{Lemma}
\newtheorem{remark}[theorem]{Remark}
\newtheorem{conjecture}[theorem]{Conjecture}

\newcommand{\E}{\mathbb E}
\newcommand{\N}{\mathbb N}
\newcommand{\R}{\mathbb R}
\newcommand{\Span}{\operatorname{span}}
\newcommand{\supp}{\operatorname{supp}}

\title{A Compactness Proof of Stable Phase Retrieval}
\author{Pedro Abdalla Teixeira\textsuperscript{1}}

\author{Jaume de Dios Pont\textsuperscript{2}}

\author{Jo\~ao Pedro Ramos\textsuperscript{3}}

\author{Mitchell Taylor\textsuperscript{4}}

\begin{document}
\begin{abstract}
We prove a stable phase retrieval result for spans of independent real-valued
random variables in $L^2$ under a uniform two-sided $L^1$ bound, allowing one
exceptional coordinate with no such hypothesis. After relabeling and
normalizing, this gives the sufficiency direction of the Calderbank,
Daubechies, Freeman, and Freeman conjecture up to one exceptional variable. The
argument is based on a compactness scheme: after the orthogonal reduction, a
counterexample sequence is decomposed into a convergent head and a
diffuse tail. The head is recovered from the quadratic observable
$|X|^2-|Y|^2$ by bounded probe functions, while the tail is controlled through
weak-limit rigidity and a uniform lower $L^1$-$L^2$ estimate for independent
sums.
\end{abstract}
\maketitle
\markboth{ABDALLA TEIXEIRA, DE DIOS PONT, RAMOS, AND TAYLOR}{A COMPACTNESS PROOF OF STABLE PHASE RETRIEVAL}
\begingroup
\renewcommand{\thefootnote}{\arabic{footnote}}
\footnotetext[1]{Department of Mathematics, University of California, Irvine, Irvine, California 92697, USA}
\footnotetext[2]{Center for Data Science, New York University, New York, New York 10011, USA}
\footnotetext[3]{Instituto de Matem\'atica Pura e Aplicada (IMPA), Rio de Janeiro, RJ, Brazil}
\footnotetext[4]{Department of Mathematics, ETH Z\"urich, Z\"urich, Switzerland}
\endgroup
\setcounter{footnote}{0}

\section{Introduction}

The phase retrieval problem is concerned with recovering a vector or function
from phaseless measurements. In the real setting, this means recovery up to
sign. The finite-dimensional theory is by now extensive; see, for instance,
\cite{bandeira2014saving,eldar2014stability} and the survey
\cite{grohs2020survey}. In infinite dimensions the situation is much more
delicate. Phase retrieval itself can hold in infinite-dimensional Hilbert
spaces \cite{cahill2016infinite}, but uniform stability fails in the general
setting of continuous frames \cite{alaifari2017continuous}, and natural
measurement models may be severely ill-posed \cite{alaifari2021gabor}.

We adopt the function-space point of view from \cite{FOTP}. If $X$ is a real
function space and $E \subset X$ is a subspace, we say that $E$ does stable
phase retrieval if there exists a constant $C < \infty$ such that
\[
\min_{\sigma \in \{\pm 1\}} \|f-\sigma g\|_X
\leq C \, \bigl\||f|-|g|\bigr\|_X
\qquad
\text{for all } f,g \in E.
\]
This formulation makes it natural to ask which subspaces of a given function
space can do stable phase retrieval. The main results of \cite{FOTP} show that
many order-continuous Banach lattices admit no infinite-dimensional closed subspaces
with this property, while \cite{christ2022examples,FOTP} provide positive
examples in $L^2$-spaces. The paper \cite{FOTP} also proves the orthogonal reduction
in real Hilbert spaces: to verify stable phase retrieval on a subspace, it is
enough to check it on orthogonal pairs. That reduction is the starting point of the present
compactness argument.

In \cite{calderbank2022stable}, Calderbank, Daubechies, Freeman, and Freeman
proved the following theorem.

\begin{theorem}
Let $(X_n)_{n=1}^\infty$ be an orthonormal sequence of independent mean-zero
random variables in $L^2([0,1])$. The following statements are equivalent.
\begin{enumerate}[label=\rm(\arabic*)]
\item The closed span of $(X_n+1_{(n,n+1]})_{n=1}^\infty$ does stable phase
retrieval in $L^2(\R)$.
\item The $L^1$ and $L^2$ norms of $(X_n)$ are uniformly comparable; that is,
there exists $A>0$ such that
\[
A\|X_n\|_{L^2([0,1])}\leq \|X_n\|_{L^1([0,1])}\leq \|X_n\|_{L^2([0,1])}
\qquad
\text{for all } n\in \N.
\]
\end{enumerate}
\end{theorem}

Following that work, the authors proposed the following conjecture.

\begin{conjecture}\label{conj:cdff}
If $(X_n)$ is an orthonormal sequence of independent mean-zero random
variables in $L^2$, then $\overline{\Span}(X_n)$ does stable phase retrieval in $L^2$ if
and only if there exist constants $0<A\leq B<1$ such that
\[
A\leq \|X_n\|_{L^1}\leq B
\]
for all but at most one $n\in \N$.
\end{conjecture}

The necessity of the conditions in Conjecture \ref{conj:cdff} is well
known.
By \cite{FOTP}, if a subspace of $L^p(\mu)$
does stable phase retrieval, then the $L^1$ and $L^p$ norms are equivalent on
that subspace, so the lower bound $A>0$ is necessary. The upper bound $B<1$
reflects the constraint that a centered random variable with
$\|f\|_{L^2}=1$ satisfies
$\|f\|_{L^1}=\|f\|_{L^2}$ if and only if it is Rademacher-valued, and the
closed span of the Rademacher functions does not do phase retrieval,
 because all Rademacher variables have equal absolute value.
Since at most one coordinate can be Rademacher, a uniform gap $B<1$ is needed
for all remaining coordinates if one requires stable phase retrieval. 

Theorem \ref{thm:main} proves the sufficiency direction with
one exceptional coordinate. After relabeling the coordinates and normalizing
each variable in $L^2$, this is exactly the case in which all but at most one
variable satisfy uniform two-sided $L^1$ bounds. Our purpose here is not to
optimize constants, but to isolate a compactness mechanism behind the proof.
After the orthogonal reduction, a counterexample sequence of pairs is split
into a convergent head and a diffuse tail whose maximal coefficient tends to
zero. The head is identified from the relation $|X_n|^2-|Y_n|^2 \to 0$ by
means of bounded probes, while the tail is controlled through a weak-limit
rigidity argument and a uniform lower $L^1$-$L^2$ estimate for independent
sums. Together these ingredients force the two coefficient vectors to agree up
to sign, contradicting the orthogonal normalization. We now turn to the precise
formulation of the theorem and the contradiction setup. The next section
reduces the proof to a normalized counterexample array; the following two
sections develop the corresponding profile decomposition and auxiliary
estimates; and the final sections identify the limiting head and tail and
derive the contradiction.

\section{Main result and contradiction setup}

We begin by stating the precise one-exception theorem proved in this note. We
then reformulate its failure in terms of a normalized orthogonal
counterexample sequence, which is the object analyzed in the rest of the
paper.

\begin{theorem}\label{thm:main}
Let $\eta,\xi_2,\xi_3,\dots$ be independent real-valued random variables and
let $\delta > 0$ be such that
\[
\|\eta\|_{L^2} = 1,
\qquad
\E[\eta] = 0,
\]
and, for every $i \ge 2$,
\[
\|\xi_i\|_{L^2} = 1,
\qquad
\E[\xi_i] = 0,
\qquad
\delta \le \|\xi_i\|_{L^1} \le 1-\delta.
\]
Then there exists $C_{\delta,\eta} \ge 1$ such that, for every
$X,Y \in \overline{\Span}(\eta,\xi_i : i \ge 2)$,
\[
\min_{\sigma \in \{\pm 1\}} \|X-\sigma Y\|_{L^2}
\le
C_{\delta,\eta} \, \bigl\| |X|-|Y| \bigr\|_{L^2}.
\]
\end{theorem}

\begin{remark}
After relabeling the coordinates and normalizing each variable in $L^2$, this
is exactly the sufficiency direction of the formulation in which
\[
\delta \le \|\xi_i\|_{L^1} \le 1-\delta
\]
holds for all but at most one index.
\end{remark}

The rest of the proof is by contradiction. After the standard orthogonality
reduction, a failure of Theorem \ref{thm:main} produces a sequence of
orthogonal pairs with almost identical moduli. The remaining sections show
that no such sequence can exist.

\begin{proposition}\label{prop:counterexample-array}
If Theorem \ref{thm:main} fails, then there exist:
\begin{enumerate}[label=\rm(\roman*)]
\item a triangular array $(\xi_{n,i})_{n \ge 1,\ i \ge 2}$ such that for each
fixed $n$ the family $(\eta,\xi_{n,i})_{i \ge 2}$ is independent and every
$\xi_{n,i}$ satisfies
\[
\|\xi_{n,i}\|_{L^2} = 1,
\qquad
\E[\xi_{n,i}] = 0,
\qquad
\delta \le \|\xi_{n,i}\|_{L^1} \le 1-\delta;
\]
\item vectors $a^{(n)}, b^{(n)} \in \ell^2(\N)$ such that
\[
\|a^{(n)}\|_2 = \|b^{(n)}\|_2 = 1,
\qquad
\langle a^{(n)}, b^{(n)} \rangle = 0,
\]
and, writing
\[
X_n := a_1^{(n)}\eta + \sum_{i \ge 2} a_i^{(n)} \xi_{n,i},
\qquad
Y_n := b_1^{(n)}\eta + \sum_{i \ge 2} b_i^{(n)} \xi_{n,i},
\]
one has
\[
\bigl\| |X_n|-|Y_n| \bigr\|_{L^2} \longrightarrow 0.
\]
\end{enumerate}
\end{proposition}

\begin{proof}
Because the variables are centered, independent, and $L^2$-normalized, each
family $(\eta,\xi_i)_{i \ge 2}$ is orthonormal in $L^2$; hence the sums
converge in $L^2$ and their $L^2$ norms equal the corresponding coefficient
$\ell^2$ norms.

If the uniform stable phase retrieval estimate fails, then for each $n$ one
may choose a violating pair $(X_n,Y_n)$. By the orthogonal reduction
\cite[Theorem~3.9]{FOTP}, this pair can be replaced by an orthogonal pair
in the same two-dimensional span without increasing the quantity
$\bigl\| |X|-|Y| \bigr\|_{L^2}$. After normalizing both functions to have
$L^2$ norm one, orthogonality gives
$\min_{\sigma=\pm 1}\|X_n-\sigma Y_n\|_{L^2}=\sqrt{2}$, so the right-hand
side must tend to zero along a suitable sequence. After relabeling the good
variables row by row, we may write
\[
X_n = a_1^{(n)}\eta + \sum_{i \ge 2} a_i^{(n)} \xi_{n,i},
\qquad
Y_n = b_1^{(n)}\eta + \sum_{i \ge 2} b_i^{(n)} \xi_{n,i},
\]
where the exceptional coordinate $\eta$ is the same for every row.
\end{proof}

\section{Profile decomposition}

Starting from the counterexample array, we next reorder the good coordinates so
that the large coefficients are moved to small indices. This yields the
head-tail decomposition that governs the rest of the argument.

\begin{lemma}[Reordered profile decomposition]\label{lem:profile}
After passing to a subsequence in $n$ and, for each row $n$, applying the same
permutation of $\{2,3,\dots\}$ to the coefficients and to the family
$(\xi_{n,i})_{i \ge 2}$, one may find vectors $a,b \in \ell^2(\N)$ and
integers $m_n \to \infty$ such that, with
\[
a^{(n)}_{\mathrm{head}} := a^{(n)} \mathbf{1}_{\{1,\dots,m_n\}},
\quad
a^{(n)}_{\mathrm{tail}} := a^{(n)} \mathbf{1}_{\{m_n+1,m_n+2,\dots\}},
\]
and likewise for $b^{(n)}$, the following hold:
\begin{enumerate}[label=\rm(\arabic*)]
\item $a^{(n)}_{\mathrm{head}} \to a$ and
$b^{(n)}_{\mathrm{head}} \to b$ strongly in $\ell^2$;
\item
\[
\max_{i > m_n} \max\bigl( |a_i^{(n)}|, |b_i^{(n)}| \bigr) \longrightarrow 0;
\]
\item
\[
\|a^{(n)}_{\mathrm{tail}}\|_2^2 \longrightarrow 1-\|a\|_2^2,
\qquad
\|b^{(n)}_{\mathrm{tail}}\|_2^2 \longrightarrow 1-\|b\|_2^2,
\]
and
\[
\langle a^{(n)}_{\mathrm{tail}}, b^{(n)}_{\mathrm{tail}} \rangle
\longrightarrow
- \langle a,b \rangle.
\]
\end{enumerate}
\end{lemma}

\begin{proof}
After passing to a subsequence, we may first assume that
\[
\bigl(a_1^{(n)},b_1^{(n)}\bigr) \longrightarrow (a_1,b_1)
\]
for some scalars $a_1,b_1 \in \R$.

For each $n$ and each $i \ge 2$, define
\[
s_i^{(n)} := \max\bigl( |a_i^{(n)}|, |b_i^{(n)}| \bigr),
\]
and permute the coordinates $\{2,3,\dots\}$ so that $(s_i^{(n)})_{i \ge 2}$ is
nonincreasing. Because
\[
\sum_{i \ge 2} \bigl(s_i^{(n)}\bigr)^2
\le
\sum_{i \ge 2} \bigl( |a_i^{(n)}|^2 + |b_i^{(n)}|^2 \bigr)
\le 2,
\]
we have the uniform bound
\[
s_i^{(n)} \le \sqrt{\frac{2}{i-1}}
\qquad
\text{for every } n \ge 1,\ i \ge 2.
\]
By a diagonal extraction we may therefore assume that, for every fixed
$i \ge 2$,
\[
a_i^{(n)} \longrightarrow a_i,
\qquad
b_i^{(n)} \longrightarrow b_i.
\]
Let
\[
a := (a_1,a_2,a_3,\dots),
\qquad
b := (b_1,b_2,b_3,\dots).
\]
Because
\[
\sum_{i \ge 1} |a_i|^2
\le
\liminf_{n \to \infty} \sum_{i \ge 1} |a_i^{(n)}|^2
= 1,
\]
and similarly for $b$, Fatou's lemma gives $a,b \in \ell^2(\N)$.

Choose an increasing sequence of integers $(m_k)_{k \ge 1}$ with $m_k \ge 1$
such that
\[
\sum_{i > m_k} |a_i|^2 + \sum_{i > m_k} |b_i|^2 \le 2^{-k}
\qquad
\text{for every } k \ge 1.
\]
For each fixed $k$, we have
\[
\sum_{i=1}^{m_k} |a_i^{(n)}-a_i|^2
+ \sum_{i=1}^{m_k} |b_i^{(n)}-b_i|^2
\longrightarrow 0
\qquad
\text{as } n \to \infty.
\]
Passing to a further subsequence in $n$, we may therefore assume that
\[
\sum_{i=1}^{m_n} |a_i^{(n)}-a_i|^2
+ \sum_{i=1}^{m_n} |b_i^{(n)}-b_i|^2
\le 2^{-n}
\qquad
\text{for every } n \ge 1.
\]
Then
\[
\|a^{(n)}_{\mathrm{head}}-a\|_2^2
\le
|a_1^{(n)}-a_1|^2
+ \sum_{i=2}^{m_n} |a_i^{(n)}-a_i|^2 + \sum_{i > m_n} |a_i|^2
\longrightarrow 0,
\]
and the same argument gives
$b^{(n)}_{\mathrm{head}} \to b$ in $\ell^2$.

Moreover, since $m_n \to \infty$ and
$s_i^{(n)} \le \sqrt{2/(i-1)}$ for $i \ge 2$, we obtain
\[
\max_{i > m_n} \max\bigl( |a_i^{(n)}|, |b_i^{(n)}| \bigr)
\le
\sqrt{\frac{2}{m_n}}
\longrightarrow 0.
\]

Finally,
\[
\|a^{(n)}_{\mathrm{tail}}\|_2^2
=
1-\|a^{(n)}_{\mathrm{head}}\|_2^2
\longrightarrow
1-\|a\|_2^2,
\]
and similarly for $b^{(n)}_{\mathrm{tail}}$. Since
\[
0 = \langle a^{(n)}, b^{(n)} \rangle
=
\langle a^{(n)}_{\mathrm{head}}, b^{(n)}_{\mathrm{head}} \rangle
+
\langle a^{(n)}_{\mathrm{tail}}, b^{(n)}_{\mathrm{tail}} \rangle,
\]
we also get
\[
\langle a^{(n)}_{\mathrm{tail}}, b^{(n)}_{\mathrm{tail}} \rangle
\longrightarrow
- \langle a,b \rangle.
\]
\end{proof}

\begin{remark}
The permutation in Lemma \ref{lem:profile} is applied only to the good
coordinates
\[
(\xi_{n,i},a_i^{(n)},b_i^{(n)})_{i \ge 2}.
\]
The exceptional coordinate $\eta$ is kept fixed at index $1$.
\end{remark}

\section{Auxiliary lemmas}

We now collect the three auxiliary ingredients used later to analyze the head
and the tail: bounded probes for recovering coefficients from the quadratic
observable, a uniform lower $L^1$-$L^2$ estimate for independent sums, and a
rigidity statement for diffuse tails.

\begin{proposition}[Bounded coefficient probes]
\label{prop:bounded-probes}
There exists $C_\delta \ge 1$ with the following property. If $\xi$ is a real
random variable satisfying
\[
\|\xi\|_{L^2}=1,
\qquad
\E[\xi] = 0,
\qquad
\delta \le \|\xi\|_{L^1} \le 1-\delta,
\]
then there exist bounded measurable functions $S,T$ such that
\[
\E[S] = \E[T] = 0,
\qquad
\E[\xi S] = 1,
\qquad
\E[(\xi^2-1)T] = 1,
\]
and
\[
\|S\|_{L^\infty}
+ \|T\|_{L^\infty}
\le C_\delta.
\]
\end{proposition}

\begin{proof}
Let
\[
S_0 := \operatorname{sign}(\xi).
\]
Then
\[
\E[\xi S_0] = \E[|\xi|] \ge \delta.
\]
Since $\E[\xi] = 0$, the centered function
\[
S := \frac{S_0 - \E[S_0]}{\E[\xi S_0]}
\]
satisfies $\E[S]=0$ and $\E[\xi S]=1$. Moreover
\[
\|S\|_{L^\infty}
\le
\frac{2}{\E[|\xi|]}
\le
\frac{2}{\delta}.
\]

Next let
\[
T_0 := \operatorname{sign}(\xi^2-1).
\]
Then
\[
\E[(\xi^2-1)T_0] = \E[|\xi^2-1|].
\]
Also,
\[
|\xi^2-1|
=
\bigl||\xi|-1\bigr| (|\xi|+1)
\ge
\bigl||\xi|-1\bigr|,
\]
so
\[
\E[|\xi^2-1|]
\ge
\E[\bigl||\xi|-1\bigr|]
\ge
\bigl| \E[|\xi|] - 1 \bigr|
=
1-\E[|\xi|]
\ge
\delta.
\]
Since $\E[\xi^2-1]=0$, the centered function
\[
T := \frac{T_0 - \E[T_0]}{\E[(\xi^2-1)T_0]}
\]
satisfies $\E[T]=0$ and $\E[(\xi^2-1)T]=1$. Finally,
\[
\|T\|_{L^\infty}
\le
\frac{2}{\E[|\xi^2-1|]}
\le
\frac{2}{\delta}.
\]
\end{proof}

\begin{remark}
For the compactness argument, Proposition \ref{prop:bounded-probes} is exactly
what is needed in order to recover the coefficients $a_i^2$ and $a_i a_j$ from
the quadratic observable $|X|^2$.
\end{remark}

\begin{proposition}[Probes for the exceptional coordinate]
\label{prop:exceptional-probes}
Let $\eta$ be a centered real-valued random variable with
\[
\|\eta\|_{L^2}=1.
\]
Then there exists a bounded measurable function $S_\eta$ such that
\[
\E[S_\eta]=0,
\qquad
\E[\eta S_\eta]=1.
\]
If $\eta^2 \not\equiv 1$ almost surely, then there exists a bounded measurable
function $T_\eta$ such that
\[
\E[T_\eta]=0,
\qquad
\E[(\eta^2-1)T_\eta]=1.
\]
\end{proposition}

\begin{proof}
Let
\[
S_0 := \operatorname{sign}(\eta).
\]
Since $\|\eta\|_{L^2}=1$, one has $\E[|\eta|]>0$. Thus
\[
S_\eta := \frac{S_0-\E[S_0]}{\E[\eta S_0]}
\]
is bounded and satisfies
\[
\E[S_\eta]=0,
\qquad
\E[\eta S_\eta]=1.
\]

Assume now that $\eta^2 \not\equiv 1$ almost surely. Let
\[
T_0 := \operatorname{sign}(\eta^2-1).
\]
Since $\E[|\,\eta^2-1\,|]>0$, the bounded function
\[
T_\eta := \frac{T_0-\E[T_0]}{\E[(\eta^2-1)T_0]}
\]
satisfies
\[
\E[T_\eta]=0,
\qquad
\E[(\eta^2-1)T_\eta]=1.
\]
\end{proof}

\begin{lemma}[Lower $L^1$-$L^2$ comparability is stable under independent sums]
\label{lem:l1l2-linear-combinations}
Let $(\zeta_i)_{i=1}^N$ be independent centered real-valued random variables and
assume that
\[
\|\zeta_i\|_{L^1} \ge \delta \|\zeta_i\|_{L^2}
\qquad
\text{for every } i=1,\dots,N.
\]
Then for every scalar family $(c_i)_{i=1}^N$ one has
\[
\left\|\sum_{i=1}^N c_i \zeta_i \right\|_{L^1}
\ge
\frac{\delta}{2\sqrt 2}
\left(\sum_{i=1}^N c_i^2 \|\zeta_i\|_{L^2}^2 \right)^{1/2}.
\]
In particular, if $\|\zeta_i\|_{L^2}=1$ for all $i$, then
\[
\left\|\sum_{i=1}^N c_i \zeta_i \right\|_{L^1}
\ge
\frac{\delta}{2\sqrt 2}\|(c_i)_{i=1}^N\|_{\ell^2}.
\]
\end{lemma}

\begin{proof}
By homogeneity, it is enough to treat the case
\[
\sum_{i=1}^N c_i^2 \|\zeta_i\|_{L^2}^2 = 1.
\]
Replacing $\zeta_i$ by $\zeta_i/\|\zeta_i\|_{L^2}$, we may therefore assume
that $\|\zeta_i\|_{L^2}=1$ and $\sum_i c_i^2=1$.

Set
\[
S := \sum_{i=1}^N c_i \zeta_i,
\]
and let
\[
S' := \sum_{i=1}^N c_i \zeta_i'
\]
be an independent copy. Then
\[
2\|S\|_{L^1}
\ge
\|S-S'\|_{L^1}
=
\left\|\sum_{i=1}^N c_i(\zeta_i-\zeta_i')\right\|_{L^1}.
\]
The differences $\zeta_i-\zeta_i'$ are independent and symmetric. Hence,
conditioning on their absolute values and applying the Khintchine inequality,
we obtain
\[
\left\|\sum_{i=1}^N c_i(\zeta_i-\zeta_i')\right\|_{L^1}
\ge
\frac{1}{\sqrt 2}
\E\left[\left(\sum_{i=1}^N c_i^2(\zeta_i-\zeta_i')^2\right)^{1/2}\right].
\]
Since $\sum_i c_i^2=1$, Cauchy--Schwarz yields the pointwise bound
\[
\sum_{i=1}^N c_i^2 |\zeta_i-\zeta_i'|
\le
\left(\sum_{i=1}^N c_i^2(\zeta_i-\zeta_i')^2\right)^{1/2},
\]
and therefore
\[
\|S\|_{L^1}
\ge
\frac{1}{2\sqrt 2}
\sum_{i=1}^N c_i^2 \|\zeta_i-\zeta_i'\|_{L^1}.
\]
Finally,
\[
\|\zeta_i-\zeta_i'\|_{L^1}
\ge
\left\|\E[\zeta_i-\zeta_i' \mid \zeta_i]\right\|_{L^1}
=
\|\zeta_i\|_{L^1}
\ge
\delta,
\]
because $\E[\zeta_i']=0$. Hence
\[
\|S\|_{L^1}
\ge
\frac{\delta}{2\sqrt 2}\sum_{i=1}^N c_i^2
=
\frac{\delta}{2\sqrt 2},
\]
which is exactly the normalized claim.
\end{proof}

\begin{proposition}[Rigidity of one-dimensional tail limits]
\label{prop:tail-rigidity}
Let $(\xi_{n,i})_{n \ge 1,\ i \ge 2}$ be a triangular array as in
Proposition \ref{prop:counterexample-array}, and let $c^{(n)} \in \ell^2(\N)$
be such that
\[
\sup_n \|c^{(n)}\|_2 < \infty
\]
and
\[
\max_{i \ge 2} |c_i^{(n)}| \longrightarrow 0,
\qquad
c_1^{(n)} = 0 \text{ for every } n.
\]
Set
\[
Z_n := \sum_{i \ge 2} c_i^{(n)} \xi_{n,i}.
\]
If $Z_n \to 0$ in distribution, then $\|c^{(n)}\|_2 \to 0$.

Equivalently, if along a subsequence the weak limit of
\[
\left(
\sum_{i \ge 2} a_i^{(n)} \xi_{n,i},
\sum_{i \ge 2} b_i^{(n)} \xi_{n,i}
\right)
\]
is supported on the line $\alpha u + \beta v = 0$ for some
$(\alpha,\beta) \ne (0,0)$, then along that same subsequence,
\[
\|\alpha a^{(n)} + \beta b^{(n)}\|_2 \longrightarrow 0.
\]
\end{proposition}

\begin{proof}
Since $Z_n \to 0$ in distribution and the limit is the constant $0$,
convergence in distribution implies $Z_n \to 0$ in probability. On the
other hand,
\[
\|Z_n\|_{L^2} = \|c^{(n)}\|_2
\]
is uniformly bounded by assumption. Thus the family $(|Z_n|)_n$ is uniformly
integrable, and therefore
\[
\|Z_n\|_{L^1} \longrightarrow 0.
\]
Applying Lemma \ref{lem:l1l2-linear-combinations} row by row to the family
$(\xi_{n,i})_{i \ge 2}$ gives
\[
\|Z_n\|_{L^1}
\ge
\frac{\delta}{2\sqrt2}\|c^{(n)}\|_2.
\]
Hence $\|c^{(n)}\|_2 \to 0$.

For the equivalent formulation, apply the first part to
\[
c_i^{(n)} := \alpha a_i^{(n)} + \beta b_i^{(n)}.
\]
The $\ell^2$ boundedness is automatic because
\[
\|c^{(n)}\|_2
\le
|\alpha|\|a^{(n)}\|_2 + |\beta|\|b^{(n)}\|_2
=
|\alpha|+|\beta|.
\]
\end{proof}

\section{Identification of the limit}

With the head-tail decomposition and the auxiliary estimates in hand, we first
identify the limiting head profile from the quadratic information
$|X_n|^2-|Y_n|^2 \to 0$, and then pass to a weak limit of the full head-tail
decomposition.

\begin{lemma}[Identification of the head profile]\label{lem:head-sign}
Either there exists $\tau \in \{\pm 1\}$ such that
\[
a = \tau b
\qquad
\text{and}
\qquad
\|a^{(n)}_{\mathrm{head}} - \tau b^{(n)}_{\mathrm{head}}\|_2 \longrightarrow 0,
\]
or else $\eta$ is a Rademacher variable and
\[
a_i=b_i=0
\qquad
\text{for every } i \ge 2.
\]
\end{lemma}

\begin{proof}
Since
\[
\bigl\| |X_n|^2 - |Y_n|^2 \bigr\|_{L^1}
\le
\|\, |X_n|-|Y_n| \,\|_{L^2}
\cdot
\|\, |X_n|+|Y_n| \,\|_{L^2}
\le
2 \|\, |X_n|-|Y_n| \,\|_{L^2}
\longrightarrow 0,
\]
fix $i \ge 2$. Apply Proposition \ref{prop:bounded-probes} to $\xi_{n,i}$ and
let $T_{n,i}$ be the corresponding probe. Note that $T_{n,i}$ is a bounded
measurable function of $\xi_{n,i}$ alone, so it is independent of
$(\eta,\xi_{n,j})_{j \ne i}$. Since $\|T_{n,i}\|_{L^\infty} \le
C_\delta$, we obtain
\[
\E\bigl[(|X_n|^2-|Y_n|^2) T_{n,i}\bigr] \longrightarrow 0.
\]
Expanding $|X_n|^2 = X_n^2$ and pairing with $T_{n,i}$, all mixed terms
vanish because the probes are centered and independent of the other
coordinates. By the defining property of $T_{n,i}$,
\[
\E\bigl[(|X_n|^2-|Y_n|^2) T_{n,i}\bigr]
=
(a_i^{(n)})^2 - (b_i^{(n)})^2,
\]
hence
\[
(a_i^{(n)})^2 - (b_i^{(n)})^2 \longrightarrow 0.
\]
Passing to the limit gives
\[
a_i^2 = b_i^2
\qquad
\text{for every } i \ge 2.
\]

Now fix $i,j \ge 2$ with $i \ne j$. Let $S_{n,i},S_{n,j}$ be the probes given by
Proposition \ref{prop:bounded-probes} for $\xi_{n,i}$ and $\xi_{n,j}$.
Again using boundedness and $L^1$ convergence,
\[
\E\bigl[(|X_n|^2-|Y_n|^2) S_{n,i} S_{n,j}\bigr] \longrightarrow 0.
\]
Expanding the square and using that $S_{n,i}$, $S_{n,j}$ are centered functions
of $\xi_{n,i}$ and $\xi_{n,j}$ respectively,
\[
\E\bigl[(|X_n|^2-|Y_n|^2) S_{n,i} S_{n,j}\bigr]
=
2\bigl(a_i^{(n)} a_j^{(n)} - b_i^{(n)} b_j^{(n)}\bigr),
\]
therefore
\[
a_i a_j = b_i b_j
\qquad
\text{for every } i,j \ge 2,\ i \ne j.
\]

Let $S_\eta$ be given by Proposition \ref{prop:exceptional-probes}. For every
$j \ge 2$, let $S_{n,j}$ be the probe given by Proposition
\ref{prop:bounded-probes} for $\xi_{n,j}$. Again using boundedness and
$L^1$ convergence,
\[
\E\bigl[(|X_n|^2-|Y_n|^2) S_\eta S_{n,j}\bigr] \longrightarrow 0.
\]
Independence gives
\[
\E\bigl[(|X_n|^2-|Y_n|^2) S_\eta S_{n,j}\bigr]
=
2\bigl(a_1^{(n)} a_j^{(n)} - b_1^{(n)} b_j^{(n)}\bigr),
\]
hence
\[
a_1 a_j = b_1 b_j
\qquad
\text{for every } j \ge 2.
\]

Assume first that there exists $i_0 \ge 2$ with $a_{i_0} \ne 0$. Since
$a_{i_0}^2=b_{i_0}^2$, we have $b_{i_0} = \tau a_{i_0}$ for some
$\tau \in \{\pm 1\}$. For every $j \ge 2$,
\[
a_{i_0} a_j = b_{i_0} b_j = \tau a_{i_0} b_j,
\]
so $a_j = \tau b_j$. Taking $j=i_0$ in the exceptional-good identity gives
\[
a_1 a_{i_0} = b_1 b_{i_0} = \tau a_{i_0} b_1,
\]
hence $a_1 = \tau b_1$. Thus $a=\tau b$.

Assume now that
\[
a_i = 0
\qquad
\text{for every } i \ge 2.
\]
Then $b_i=0$ for every $i \ge 2$ as well. If $\eta^2 \not\equiv 1$ almost
surely, let $T_\eta$ be given by Proposition
\ref{prop:exceptional-probes}. Again using boundedness and $L^1$ convergence,
\[
\E\bigl[(|X_n|^2-|Y_n|^2) T_\eta\bigr] \longrightarrow 0.
\]
By independence,
\[
\E\bigl[(|X_n|^2-|Y_n|^2) T_\eta\bigr]
=
(a_1^{(n)})^2-(b_1^{(n)})^2,
\]
so $a_1^2=b_1^2$ and therefore $a=\tau b$ for some $\tau \in \{\pm 1\}$.

Since $\E[\eta]=0$ and $\|\eta\|_{L^2}=1$, the remaining case
$\eta^2 \equiv 1$ almost surely means that $\eta$ is a Rademacher variable.
This is exactly the second alternative.

If we are in the first alternative, then
\[
\|a^{(n)}_{\mathrm{head}} - \tau b^{(n)}_{\mathrm{head}}\|_2
\le
\|a^{(n)}_{\mathrm{head}}-a\|_2
+
\|b^{(n)}_{\mathrm{head}}-b\|_2
\longrightarrow 0.
\]
\end{proof}

\begin{lemma}[Limit law for the head-tail decomposition]
\label{lem:limit-law}
After passing to a further subsequence, the quadruple
\[
\bigl(
X_n^{\mathrm{head}},
Y_n^{\mathrm{head}},
X_n^{\mathrm{tail}},
Y_n^{\mathrm{tail}}
\bigr)
\]
converges in distribution to a random vector
\[
\bigl(
H_X,H_Y,R_X,R_Y
\bigr)
\in \R^4,
\]
where:
\begin{enumerate}[label=\rm(\arabic*)]
\item $(H_X,H_Y)$ is independent of $(R_X,R_Y)$;
\item $(R_X,R_Y)$ is infinitely divisible;
\item either $H_X = \tau H_Y$ almost surely for some $\tau \in \{\pm 1\}$, or
else $\eta$ is a Rademacher variable and
\[
H_X = a_1 \eta,
\qquad
H_Y = b_1 \eta
\qquad
\text{almost surely;}
\]
\item
\[
|H_X + R_X| = |H_Y + R_Y|
\qquad
\text{almost surely.}
\]
\end{enumerate}
\end{lemma}

\begin{proof}
Set
\[
X_n^{\mathrm{head}} := a_1^{(n)}\eta + \sum_{2 \le i \le m_n} a_i^{(n)} \xi_{n,i},
\qquad
X_n^{\mathrm{tail}} := \sum_{i > m_n} a_i^{(n)} \xi_{n,i},
\]
and define $Y_n^{\mathrm{head}}$, $Y_n^{\mathrm{tail}}$ analogously.
Each of these random variables is bounded in $L^2$ by $1$, so the family of
quadruples is tight. After passing to a subsequence we therefore get a weak
limit $(H_X,H_Y,R_X,R_Y)$.

For each $n$, the pair
$(X_n^{\mathrm{head}},Y_n^{\mathrm{head}})$ depends only on $\eta$ and on
$(\xi_{n,i})_{2 \le i \le m_n}$, while
$(X_n^{\mathrm{tail}},Y_n^{\mathrm{tail}})$ depends only on
$(\xi_{n,i})_{i > m_n}$. The two pairs are independent. The factorization of
characteristic functions is preserved under weak convergence, hence
$(H_X,H_Y)$ is independent of $(R_X,R_Y)$.

To see that $(R_X,R_Y)$ is infinitely divisible, consider the row-wise array of
summands
\[
Z_{n,i}
:=
\bigl(a_i^{(n)} \xi_{n,i},\, b_i^{(n)} \xi_{n,i}\bigr)\mathbf{1}_{\{i>m_n\}}.
\]
The array is infinitesimal because, for every $\varepsilon>0$,
\[
\sup_{i>m_n}
\mathbb P\bigl(|Z_{n,i}|>\varepsilon\bigr)
\le
\frac{1}{\varepsilon^2}
\sup_{i>m_n}\bigl((a_i^{(n)})^2+(b_i^{(n)})^2\bigr)
\longrightarrow 0.
\]
For each fixed $n$, the summands $(Z_{n,i})_{i > m_n}$ are independent since
the $\xi_{n,i}$ are independent. By the classical theorem on limits of
row-wise independent infinitesimal triangular arrays
\cite[Chapter~4]{gnedenko-kolmogorov}, every subsequential weak limit of the
sums $\sum_i Z_{n,i}$ is infinitely divisible.
This gives (2).

If we are in the first alternative of Lemma \ref{lem:head-sign}, then
\[
\|X_n^{\mathrm{head}} - \tau Y_n^{\mathrm{head}}\|_{L^2}
=
\|a^{(n)}_{\mathrm{head}} - \tau b^{(n)}_{\mathrm{head}}\|_2
\longrightarrow 0,
\]
hence $H_X = \tau H_Y$ almost surely.

If we are in the second alternative of Lemma \ref{lem:head-sign}, then
\[
\sum_{2 \le i \le m_n} (a_i^{(n)})^2
\longrightarrow 0,
\qquad
\sum_{2 \le i \le m_n} (b_i^{(n)})^2
\longrightarrow 0,
\]
and therefore
\[
\|X_n^{\mathrm{head}} - a_1 \eta\|_{L^2}
\le
|a_1^{(n)}-a_1|
+
\left(\sum_{2 \le i \le m_n} (a_i^{(n)})^2\right)^{1/2}
\longrightarrow 0,
\]
with the same argument for $Y_n^{\mathrm{head}}$. Thus
\[
H_X = a_1 \eta,
\qquad
H_Y = b_1 \eta
\qquad
\text{almost surely.}
\]

Finally,
\[
\bigl|X_n^{\mathrm{head}} + X_n^{\mathrm{tail}}\bigr|
-
\bigl|Y_n^{\mathrm{head}} + Y_n^{\mathrm{tail}}\bigr|
=
|X_n|-|Y_n|
\longrightarrow 0
\]
in $L^2$, hence in probability. On the other hand, by the continuous mapping
theorem this difference converges in distribution to
$|H_X+R_X|-|H_Y+R_Y|$. Therefore the limit must vanish almost surely.
\end{proof}

\begin{lemma}[A geometric two-line support lemma]
\label{lem:two-line}
Let $\mu$ be an infinitely divisible probability measure on $\R^2$. If
$\supp \mu$ is contained in the union of two affine lines, then
$\supp \mu$ is contained in a single affine line.
\end{lemma}

\begin{proof}
Because $\mu$ is infinitely divisible, it is in particular $2$-divisible:
there exists a probability measure $\nu$ such that $\mu=\nu * \nu$.
Let $A := \supp \nu$. Then
\[
\supp \mu = \overline{A+A}.
\]

Assume for contradiction that $A$ is not contained in an affine line. Choose
three non-collinear points $x,y,z \in A$. Let $L_1 \cup L_2$ be the union of
two affine lines containing $\supp \mu$. Since $A+A \subset \supp \mu$, the six
points
\[
2x,\ 2y,\ 2z,\ x+y,\ x+z,\ y+z
\]
belong to $L_1 \cup L_2$. Among the three points $2x,2y,2z$, two must belong
to the same line; after relabeling, assume $2x,2y \in L_1$.

Then $L_1$ is the line through $2x$ and $2y$. If $x+z \in L_1$, then $x,y,z$
would be collinear, which is impossible. The same argument shows that
$y+z \notin L_1$ and $2z \notin L_1$. Consequently,
\[
2z,\ x+z,\ y+z \in L_2.
\]
But $x+z$ and $y+z$ lie on the line $z+\R(y-x)$, whereas $2z$ belongs to that
line only if $z,x,y$ are collinear. This contradiction proves that $A$ is
contained in an affine line, and therefore so is
$\supp \mu = \overline{A+A}$.
\end{proof}

\section{Proof of Theorem \ref{thm:main}}

We can now complete the contradiction argument by combining the structural
information obtained for the head with the rigidity of the tail limit.

\begin{proof}
Assume that Theorem \ref{thm:main} fails, and let the counterexample array be
given by Proposition \ref{prop:counterexample-array}. Apply
Lemma \ref{lem:profile}, then Lemma \ref{lem:head-sign}, and finally
Lemma \ref{lem:limit-law}. We use the notation introduced there.

Since $(H_X,H_Y)$ is independent of $(R_X,R_Y)$, we may realize the joint law
on a product space. By Fubini's theorem, there exists a set $\Omega_0$ of full
probability for the head variables such that, for every $\omega \in \Omega_0$,
\[
|H_X(\omega)+R_X| = |H_Y(\omega)+R_Y|
\qquad
\text{holds almost surely in the tail variables.}
\]
Assume first that we are in the second alternative of Lemma
\ref{lem:head-sign}. Then $\eta$ is a Rademacher variable and, by Lemma
\ref{lem:limit-law},
\[
H_X = a_1 \eta,
\qquad
H_Y = b_1 \eta
\qquad
\text{almost surely.}
\]
Choose $\omega_+,\omega_- \in \Omega_0$ such that
\[
\bigl(H_X(\omega_+),H_Y(\omega_+)\bigr)=(a_1,b_1),
\qquad
\bigl(H_X(\omega_-),H_Y(\omega_-)\bigr)=(-a_1,-b_1).
\]
Then the support of the law of $(R_X,R_Y)$ is contained in both
\[
\Gamma_+
:=
\{(u,v)\in\R^2 : u-v=b_1-a_1\}
\cup
\{(u,v)\in\R^2 : u+v=-(a_1+b_1)\}
\]
and
\[
\Gamma_-
:=
\{(u,v)\in\R^2 : u-v=a_1-b_1\}
\cup
\{(u,v)\in\R^2 : u+v=a_1+b_1\}.
\]
By Lemma \ref{lem:two-line}, the support of the law of $(R_X,R_Y)$ is
contained in a single affine line. Since $(R_X,R_Y)$ is centered (as a weak
limit of centered random vectors), the barycenter of the law lies at the
origin, so the line must pass through the origin.

If $(R_X,R_Y)=(0,0)$ almost surely, then Proposition
\ref{prop:tail-rigidity} applied separately to
\[
\sum_{i > m_n} a_i^{(n)} \xi_{n,i}
\qquad
\text{and}
\qquad
\sum_{i > m_n} b_i^{(n)} \xi_{n,i}
\]
gives
\[
\|a^{(n)}_{\mathrm{tail}}\|_2 + \|b^{(n)}_{\mathrm{tail}}\|_2 \longrightarrow 0.
\]
Since $a_i=b_i=0$ for every $i \ge 2$, Lemma \ref{lem:profile} yields
\[
\|a^{(n)}_{\mathrm{head}}\|_2 \longrightarrow |a_1|,
\qquad
\|b^{(n)}_{\mathrm{head}}\|_2 \longrightarrow |b_1|.
\]
The norm identities therefore force $|a_1|=|b_1|=1$. On the other hand,
\[
0
=
\langle a^{(n)},b^{(n)}\rangle
=
\langle a^{(n)}_{\mathrm{head}},b^{(n)}_{\mathrm{head}}\rangle
+
\langle a^{(n)}_{\mathrm{tail}},b^{(n)}_{\mathrm{tail}}\rangle,
\]
and the second term tends to $0$. Passing to the limit gives
$a_1 b_1 = 0$, contradicting $|a_1|=|b_1|=1$.

Otherwise choose a nonzero point in the support of $(R_X,R_Y)$. Since the
support is contained in a line through the origin and in both $\Gamma_+$ and
$\Gamma_-$, that line must be one of $u=v$ or $u=-v$.

If the support is contained in $u=v$, choose a nonzero point $(u,u)$ in the
support. Since $(u,u) \in \Gamma_+ \cap \Gamma_-$, one must have $a_1=b_1$.
If the support is contained in $u=-v$, choose a nonzero point $(u,-u)$ in the
support. Since $(u,-u) \in \Gamma_+ \cap \Gamma_-$, one must have $a_1=-b_1$.

Thus there exists $\sigma \in \{\pm 1\}$ such that
\[
R_X=\sigma R_Y \text{ almost surely}
\qquad
\text{and}
\qquad
a_1=\sigma b_1.
\]
Proposition \ref{prop:tail-rigidity} yields
\[
\|a^{(n)}_{\mathrm{tail}} - \sigma b^{(n)}_{\mathrm{tail}}\|_2 \longrightarrow 0.
\]
Since $a=(a_1,0,0,\dots)$ and $b=(b_1,0,0,\dots)$ with $a_1=\sigma b_1$,
Lemma \ref{lem:profile} gives
\[
\|a^{(n)}_{\mathrm{head}} - \sigma b^{(n)}_{\mathrm{head}}\|_2 \longrightarrow 0.
\]
Hence
\[
\|a^{(n)}-\sigma b^{(n)}\|_2 \longrightarrow 0,
\]
again contradicting
\[
\|a^{(n)}-\sigma b^{(n)}\|_2^2 = 2.
\]

We are left with the first alternative of Lemma \ref{lem:head-sign}. Fix
$\tau \in \{\pm 1\}$ such that
\[
H_X=\tau H_Y \text{ almost surely.}
\]
Then for every $\omega \in \Omega_0$ the support of the law of $(R_X,R_Y)$ is
contained in
\[
\Gamma_\omega :=
\begin{cases}
\{(u,v)\in\R^2 : u-v=0\} \cup \{(u,v)\in\R^2 : u+v=-2H_X(\omega)\},
& \tau=1, \\[2mm]
\{(u,v)\in\R^2 : u+v=0\} \cup \{(u,v)\in\R^2 : u-v=-2H_X(\omega)\},
& \tau=-1.
\end{cases}
\]
By Lemma \ref{lem:two-line}, the support of the law of $(R_X,R_Y)$ is
contained in a single affine line. Since $(R_X,R_Y)$ is centered, that line
must pass through the origin.

We now distinguish two cases.

\medskip
\noindent
\textbf{Case 1: $H_X \not\equiv 0$.}
Choose $\omega \in \Omega_0$ with $H_X(\omega)\ne 0$. Then among the two lines
forming $\Gamma_\omega$, the only one passing through the origin is
\[
\{u-v=0\} \quad \text{if } \tau=1,
\qquad
\{u+v=0\} \quad \text{if } \tau=-1.
\]
Therefore the law of $(R_X,R_Y)$ is supported on the line $u=\tau v$. By
Proposition \ref{prop:tail-rigidity},
\[
\|a^{(n)}_{\mathrm{tail}} - \tau b^{(n)}_{\mathrm{tail}}\|_2 \longrightarrow 0.
\]
Combining this with
\[
\|a^{(n)}_{\mathrm{head}} - \tau b^{(n)}_{\mathrm{head}}\|_2 \longrightarrow 0
\]
from Lemma \ref{lem:head-sign}, we obtain
\[
\|a^{(n)}-\tau b^{(n)}\|_2 \longrightarrow 0.
\]
This is impossible because
\[
\|a^{(n)}-\tau b^{(n)}\|_2^2
=
\|a^{(n)}\|_2^2 + \|b^{(n)}\|_2^2 - 2\tau \langle a^{(n)},b^{(n)}\rangle
= 2.
\]

\medskip
\noindent
\textbf{Case 2: $H_X \equiv 0$.}
Since $X_n^{\mathrm{head}} \to H_X = 0$ in distribution and
$\|X_n^{\mathrm{head}}\|_{L^2} \le 1$, uniform integrability gives
$\|X_n^{\mathrm{head}}\|_{L^1} \to 0$.

Let $S_\eta$ be given by Proposition \ref{prop:exceptional-probes}. Then
\[
a_1^{(n)}
=
\E[X_n^{\mathrm{head}} S_\eta]
\longrightarrow 0,
\]
because $\|S_\eta\|_{L^\infty} < \infty$ and
$|\E[X_n^{\mathrm{head}} S_\eta]| \le \|S_\eta\|_{L^\infty}
\|X_n^{\mathrm{head}}\|_{L^1} \to 0$.
For every fixed $i \ge 2$, one has $i \le m_n$ for all sufficiently large $n$.
Let $S_{n,i}$ be the probe given by Proposition \ref{prop:bounded-probes} for
$\xi_{n,i}$. Then
\[
a_i^{(n)}
=
\E[X_n^{\mathrm{head}} S_{n,i}]
\longrightarrow 0,
\]
because $\|S_{n,i}\|_{L^\infty} \le C_\delta$ and
$|\E[X_n^{\mathrm{head}} S_{n,i}]| \le C_\delta
\|X_n^{\mathrm{head}}\|_{L^1} \to 0$. Passing to
the limit gives $a_i=0$ for every $i \ge 1$, hence $a=0$.

The same argument gives $b=0$, so Lemma \ref{lem:profile} yields
\[
\|a^{(n)}_{\mathrm{head}}\|_2 + \|b^{(n)}_{\mathrm{head}}\|_2 \longrightarrow 0.
\]
Since $H_X = \tau H_Y = 0$, the set $\Gamma_\omega$ reduces to
$\{u=v\} \cup \{u=-v\}$ for every $\omega$. By Lemma \ref{lem:two-line}, the
support of the law of $(R_X,R_Y)$ is contained in one of the two lines $u=v$
or $u=-v$. Thus there exists $\sigma \in \{\pm 1\}$ such that
$R_X=\sigma R_Y$ almost surely, and Proposition
\ref{prop:tail-rigidity} yields
\[
\|a^{(n)}_{\mathrm{tail}} - \sigma b^{(n)}_{\mathrm{tail}}\|_2 \longrightarrow 0.
\]
Therefore
\[
\|a^{(n)}-\sigma b^{(n)}\|_2
\le
\|a^{(n)}_{\mathrm{head}}\|_2
+
\|b^{(n)}_{\mathrm{head}}\|_2
+
\|a^{(n)}_{\mathrm{tail}}-\sigma b^{(n)}_{\mathrm{tail}}\|_2
\longrightarrow 0,
\]
again contradicting
\[
\|a^{(n)}-\sigma b^{(n)}\|_2^2 = 2.
\]

The contradiction completes the proof.
\end{proof}

\bibliographystyle{plain}
\bibliography{refs.bib}

\end{document}
